% =====================================================================
%  PREÁMBULO COMÚN (memoria, informes y guía)
%  Normalmente NO hace falta tocar este fichero: los datos van en
%  config/datos.tex. Si añades paquetes, hazlo al final, en la zona
%  marcada, y comprueba que el documento sigue compilando sin errores.
%  Requiere: \idiomadocumento y \estilocitas (definidos en datos.tex).
% =====================================================================

% --- Codificación, fuentes y microtipografía --------------------------
\usepackage[T1]{fontenc}
\usepackage[utf8]{inputenc}
\usepackage{lmodern}
\usepackage[babel=true]{microtype}

% --- Idiomas: castellano, catalán e inglés; el principal sale de datos.tex
%     es-tabla        -> «Tabla» en vez de «Cuadro»
%     es-noshorthands -> evita conflictos de < > " con TikZ y tablas
%     es-nodecimaldot -> 3.5 en fórmulas no se convierte en 3,5
\edef\tfgbabelopts{english,catalan,spanish,main=\idiomadocumento}
\expandafter\usepackage\expandafter[\tfgbabelopts,es-tabla,es-noshorthands,es-nodecimaldot]{babel}
\usepackage{csquotes}

% --- Página: A4, márgenes, interlineado 1,15, párrafos sin línea en blanco
\usepackage[a4paper,top=2.5cm,bottom=2.5cm,left=3cm,right=2.5cm,headheight=15pt]{geometry}
\usepackage{setspace}
\setstretch{1.15}
\usepackage[skip=6pt plus 2pt]{parskip}   % separación entre párrafos sin dobles saltos
\raggedbottom
% Sin líneas viudas ni huérfanas
\clubpenalty=10000
\widowpenalty=10000
\displaywidowpenalty=10000

% --- Colores, figuras, tablas y diagramas -----------------------------
\usepackage[table,dvipsnames]{xcolor}
\usepackage{graphicx}
\graphicspath{{figuras/}}
\usepackage{float}
\usepackage[section]{placeins}            % las figuras no "huyen" a otra sección
\usepackage{booktabs,tabularx,multirow,array,longtable}
\newcolumntype{L}{>{\raggedright\arraybackslash}X}
\newcolumntype{C}{>{\centering\arraybackslash}X}
\renewcommand{\arraystretch}{1.2}
\usepackage{amsmath,amssymb}
\usepackage{siunitx}                      % \qty{100}{ms}, \num{12345}: cifras y unidades
\usepackage{tikz}
\usetikzlibrary{arrows.meta,positioning,fit,backgrounds,shapes.geometric,babel}
\usepackage{pgfgantt}
\usepackage{qrcode}
\usepackage{enumitem}
\setlist{itemsep=2pt,topsep=4pt}

% --- Pies de figura y tabla: centrados, breves; tablas con el pie ARRIBA
\usepackage[font=small,labelfont=bf,labelsep=colon,justification=centering,
            singlelinecheck=true]{caption}
\captionsetup[table]{position=top,skip=6pt}
\captionsetup[figure]{skip=6pt}
\usepackage[list=true]{subcaption}

% \fuente{...}: línea de fuente bajo la figura o la tabla (OBLIGATORIA).
%   \fuente{Elaboración propia.}   \fuente{Adaptado de \cite{clave}.}
\newcommand{\fuente}[1]{\par\vspace{3pt}{\centering\footnotesize\fuentename: #1\par}}

% --- Código fuente --------------------------------------------------
\usepackage{listings}
\lstset{
  basicstyle=\ttfamily\small, columns=fullflexible, keepspaces=true,
  breaklines=true, breakatwhitespace=false, tabsize=2,
  frame=tb, framesep=4pt, numbers=left, numberstyle=\tiny\color{gray}, numbersep=6pt,
  keywordstyle=\color{RoyalBlue}\bfseries, commentstyle=\color{ForestGreen}\itshape,
  stringstyle=\color{BrickRed}, showstringspaces=false, captionpos=b,
  floatplacement=htbp,  % con [float] en cada listado, el código no se parte entre páginas
  extendedchars=true,
  literate={á}{{\'a}}1 {é}{{\'e}}1 {í}{{\'i}}1 {ó}{{\'o}}1 {ú}{{\'u}}1
           {Á}{{\'A}}1 {É}{{\'E}}1 {Í}{{\'I}}1 {Ó}{{\'O}}1 {Ú}{{\'U}}1
           {à}{{\`a}}1 {è}{{\`e}}1 {ò}{{\`o}}1 {ï}{{\"i}}1 {ü}{{\"u}}1
           {ñ}{{\~n}}1 {Ñ}{{\~N}}1 {ç}{{\c c}}1 {Ç}{{\c C}}1 {·}{{$\cdot$}}1
           {¿}{{?`}}1 {¡}{{!`}}1
}

% --- Acrónimos: \ac{TFG} -> 1.ª vez «Trabajo de Fin de Grado (TFG)», después «TFG»
\usepackage[printonlyused]{acronym}

% --- Bibliografía (biblatex + biber). Un ÚNICO estilo en todo el documento.
\edef\tfgbibopts{backend=biber,style=\estilocitas}
\expandafter\usepackage\expandafter[\tfgbibopts]{biblatex}
\usepackage{xurl}                          % URLs largas que se parten bien
% Fecha de consulta y URL en el idioma del documento (nunca «Accessed» en un
% texto en castellano):  «Consultado el 15 de nov. de 2026. Disponible en: URL»
\DefineBibliographyStrings{spanish}{urlseen={Consultado el}, url={Disponible en}}
\DefineBibliographyStrings{catalan}{urlseen={Consultat el}, url={Disponible a}}

% --- Textos que dependen del idioma --------------------------------------
\addto\captionsspanish{%
  \def\fuentename{Fuente}\def\acronimosname{Acrónimos}%
  \def\lstlistingname{Listado}\def\lstlistlistingname{Índice de listados}%
  \def\resumenname{Resumen}\def\clavesname{Palabras clave}%
  \sisetup{output-decimal-marker={,}}}
\addto\captionscatalan{%
  \def\fuentename{Font}\def\acronimosname{Acrònims}%
  \def\lstlistingname{Llistat}\def\lstlistlistingname{Índex de llistats}%
  \def\resumenname{Resum}\def\clavesname{Paraules clau}%
  \sisetup{output-decimal-marker={,}}}
\addto\captionsenglish{%
  \def\fuentename{Source}\def\acronimosname{Acronyms}%
  \def\lstlistingname{Listing}\def\lstlistlistingname{List of Listings}%
  \def\resumenname{Abstract}\def\clavesname{Keywords}%
  \sisetup{output-decimal-marker={.}}}

% --- Encabezados de capítulo y profundidad del índice ----------------
\usepackage{titlesec}
\ifdefined\chapter
  \titleformat{\chapter}[display]
    {\normalfont\bfseries\huge\raggedright\hyphenpenalty=10000}
    {\large\chaptertitlename~\thechapter}{6pt}{\huge}
    [\vspace{2pt}\titlerule]
  \titlespacing*{\chapter}{0pt}{0pt}{20pt}
\fi
% Títulos sin palabras partidas por guiones
\titleformat*{\section}{\normalfont\Large\bfseries\raggedright\hyphenpenalty=10000}
\titleformat*{\subsection}{\normalfont\large\bfseries\raggedright\hyphenpenalty=10000}
\setcounter{secnumdepth}{2}   % numera hasta X.Y.Z; no uses un 4.º nivel
\setcounter{tocdepth}{1}      % el índice muestra capítulos y secciones (máx. 2 niveles)

% --- Cabecera y pie: número de página centrado abajo ----------------
\usepackage{fancyhdr}
\pagestyle{fancy}
\fancyhf{}
\fancyhead[R]{\small\nouppercase{\leftmark}}
\fancyfoot[C]{\thepage}
\renewcommand{\headrulewidth}{0.4pt}
\fancypagestyle{plain}{\fancyhf{}\fancyfoot[C]{\thepage}\renewcommand{\headrulewidth}{0pt}}

% --- Enlaces y metadatos del PDF (hyperref casi al final) --------------
\usepackage[colorlinks=true,linkcolor=black,citecolor=black,
            urlcolor=MidnightBlue,bookmarksnumbered=true,pdfusetitle]{hyperref}
\usepackage{bookmark}

% --- Zona para tus paquetes y comandos propios -----------------------
% \usepackage{...}
